%---------------------------------------------------------------------------%
%->> 深圳大学毕业论文模板
%---------------------------------------------------------------------------%
%- 载入模板类
\documentclass{szuthesis}% 默认形式
% \documentclass[print]{szuthesis}% 打印预览版本，可以自动生成额外的空白页用于打印
% \documentclass[fontset=windows|adobe|mac|ubuntu]{szuthesis}% 选择字库
%---------------------------------------------------------------------------%
% - 载入配置信息，包含论文封面信息、必要的package
%---------------------------------------------------------------------------%
%->> Cover information 封面信息
%---------------------------------------------------------------------------%
\classid{TP391}% 分类号
\udc{004}
\confidential{公开}% 密级
%- 注：\title包含两个参数
% \title{深圳大学\LaTeX{}模板}{}% 单行题目，第二个参数为空
\title{新巴塞尔协议风险管理理念与}{我国风险管理体系的构建}% 多行题目
%- 注：英文题目用于生成Abstract的页眉，只有一个参数
\TITLE{The Idea of Global Risk Management from the New Basle Capital Accord and the ERM's Constructing of in China}
\author{唐同学}% 论文作者
\major{统计学}% 学科专业名称
%- 注：以下两个类型支持多行
\institute{计算机与软件学院}% 院系名称单行
% \institute{某某学院\\某某实验室}% 院系名称多行
\advisor{张老师\ 教授}% 指导教师单行
% \advisor{张老师\ 教授\\王老师\ 研究员}% 指导教师多行
\DEGREE{MasterXS}% 学术硕士
% \DEGREE{MasterZY}\degree{此处填写学位类别}% 专业硕士
%---------------------------------------------------------------------------%
%->> other config
%---------------------------------------------------------------------------%
%- 添加两个命令方便输出
\DeclareRobustCommand\cs[1]{\texttt{\char`\\#1}}
\providecommand\pkg[1]{{\sffamily#1}}
%-
\addbibresource{Biblio/ref.bib}% 参考文献路径
\setlength\bibitemsep{0.0ex plus 0.2ex minus 0.2ex}% set distance between bib entrie
%-
\setcounter{tocdepth}{3}% depth for the table of contents，设为2可不显示subsubsection
\setcounter{secnumdepth}{3}% depth for section numbering, default is 2
%-
%- 某些小语种会超出版面边界，提示Overfull \hbox{}...，中英日韩无需使用(或使用宏包microtype)
% \setlength\emergencystretch{1em}
%-
%- 重新设置 equation, figure, table 的序号
%\numberwithin{equation}{section}% set enumeration level
%\renewcommand{\theequation}{\thesection\arabic{equation}}% configure the label style
%\numberwithin{figure}{section}% set enumeration level
%\renewcommand{\thefigure}{\thesection\arabic{figure}}% configure the label style
%\numberwithin{table}{section}% set enumeration level
%\renewcommand{\thetable}{\thesection\arabic{table}}% configure the label style
\counterwithout{footnote}{chapter}% footnote编号全局连续
%-
%---------------------------------------------------------------------------%
%->> Package
%---------------------------------------------------------------------------%
% -> szuthesis.cls中已经导入的包
% - etoolbox, a toolbox of programming facilities
% - geometry, for layout
% - expl3, LaTeX3 programming environment
% - array
% - ulem, underline
% - xeCJKfntef, underline for CJK
% - fancyhdr, header and footer
% - biblatex
%-
\usepackage{graphicx}
\DeclareGraphicsExtensions{.pdf,.jpg,.png,.eps,.tif,.bmp}% 默认图片格式
\graphicspath{{Image/}}% 默认图片检索路径
%-
\usepackage[format=plain,hangindent=2.0em,font={small},skip=8pt,labelsep=space]{caption}
%-
\usepackage{subcaption}% 处理子图
%-
% \usepackage[list=off]{bicaption} % 双语caption
% \DeclareCaptionOption{bi-second}[]{
%     \def\tablename{Table}%
%     \def\figurename{Figure}%
% }
% \captionsetup[bi-second]{bi-second}
%-
\usepackage[section]{placeins}% 阻止图片浮动超出当前section
%-
\usepackage{enumitem}% 列表环境功能提升
\setlist{nosep}% 默认文本行间距
% \setlist[enumerate]{wide=\parindent}% 是否悬挂对齐，不建议全局修改
% \setlist[itemize]{wide=\parindent}
%-
% \usepackage{verbatim}
%-
% \usepackage{chemfig}% draw 2D chemical structures
% \usepackage[version=4]{mhchem}% typeset chemical formulae [mhchem|chemformula]
%-
% \usepackage{microtype}% improves general appearance of the text, 启用后降低编译效率
%-
% \usepackage{pdflscape}% landscape environment, \begin{landscape} ... \end{landscape}
%-
% \usepackage[usenames,dvipsnames,svgnames,table]{xcolor}% color support
%-
% \usepackage{tikz}% automatically load pgf package, plot with tex
% \usetikzlibrary{positioning, arrows, calc, trees }%
%-
\usepackage{booktabs}% 三线表
%-
\usepackage{listings}% 代码片段
\def\lstlistingname{代码}
\lstset{%
    basicstyle=\linespread{1.2}\small, % 字体
    breaklines=true,                   % 自动换行
    frame=lines,                       % 上下的边框，可选none|single|shadowbox等
    keepspaces=true,
    showstringspaces=false,            % string的空格添加标记，defaul:true
    tabsize=2,                         % tab长度
    % stringstyle=\color{DarkViolet},
    % backgroundcolor=\color{gray!10},
    % commentstyle=\color{ForestGreen},
    % keywordstyle=\color{blue},
}
%-
\usepackage[ruled,vlined,linesnumbered,algochapter]{algorithm2e} % 算法描述
\SetAlgorithmName{算法}{算法}{}
%---------------------------------------------------------------------------%
%->> 配置数学环境
%---------------------------------------------------------------------------%
\usepackage{amsmath}
\usepackage{amssymb}
%- 符号表，参考 http://milde.users.sourceforge.net/LUCR/Math/mathpackages/amssymb-symbols.pdf
\usepackage{amsthm} % 定理引理等环境
\theoremstyle{plain}% for theorems, lemmas, propositions, etc
\newtheorem{theorem}              {定理} [chapter]
\newtheorem{axiom}      [theorem] {公理}
\newtheorem{lemma}      [theorem] {引理}
\newtheorem{corollary}  [theorem] {推论}
\newtheorem{assertion}  [theorem] {断言}
\newtheorem{proposition}[theorem] {命题}
\newtheorem{conjecture} [theorem] {猜想}
\newtheorem{assumption} [theorem] {假设}
\theoremstyle{definition}% for definitions and examples
\newtheorem{definition}           {定义} [chapter]
\newtheorem{example}              {例}   [chapter]
\newtheorem{problem}              {问题} [chapter]
\newtheorem{exercise}             {练习} [chapter]
\theoremstyle{remark}% for remarks and notes
\newtheorem*{remark}              {注}
\newtheorem*{solution}            {解}
% \usepackage{mathtools}
\usepackage{unicode-math}
%- 注：unicode-math可以配置数学公式字体，注意包冲突！
%- 已知可能存在冲突的包：amscd,amsfonts,bbm,bm,eucal,eufrak,mathrsfs
\setmathfont{XITSMath-Regular}[
    Extension=.otf, BoldFont=XITSMath-Bold, Ligatures=TeX, StylisticSet = 1,
]
\setmathfont{XITSMath-Regular}[
    Extension=.otf, range={scr,bfscr}, Ligatures=TeX, StylisticSet = 2,
]
\setmathfont{XITSMath-Regular}[
    Extension=.otf, range={cal,bfcal}, Ligatures=TeX, StylisticSet = 1,
]
% \setmathfont{XITS Math Bold}[version=bold]% for bold version % 不兼容StylisticSet=2
% \newenvironment{szumathbf}{\bfseries\mathversion{bold}}{}
\def\XITSMathFontOptions{
    Extension=.otf, BoldFont=XITSMath-Bold, Ligatures=TeX, StylisticSet = 1
}
\setmathrm{XITSMath-Regular}[\XITSMathFontOptions]
\setmathsf{XITSMath-Regular}[\XITSMathFontOptions]
\setmathtt{XITSMath-Regular}[\XITSMathFontOptions]
%-
\def\boldsymbol#1{\symbfit{#1}}
\providecommand{\Vector}[1]{\symbfit{#1}}
\providecommand{\Matrix}[1]{\symbfit{#1}}
\providecommand{\Tensor}[1]{\symbfit{#1}}
\providecommand{\Dif}{\symrm{d}}
\providecommand{\Const}[1]{\symrm{#1}}
\providecommand{\deltarm}{\symrm{\delta}}
\providecommand{\Div}{\operatorname{div}}
\providecommand{\Trace}{\operatorname{tr}}
%---------------------------------------------------------------------------%
%->> 链接，生成书签，在最后
%---------------------------------------------------------------------------%
\usepackage{hyperref}% 超链接，生成书签，[注：放在最后]
\hypersetup{% set hyperlinks
    pdfencoding=auto,% allows non-Latin based languages in bookmarks
    psdextra=true,% extra support for math symbols in bookmarks
    bookmarksnumbered=true,% put section numbers in bookmarks
    pdftitle={\szutitle},% title
    pdfauthor={\szuauthor},% author
    pdfsubject={\szutitle},% subject
    pdfstartview={FitH},% fits the width of the page to the window
    % colorlinks=true,% false: boxed links; true: colored links
    % linkcolor=black,% color of internal links
    % citecolor=blue,% color of links to bibliography
    % filecolor=blue,% color of file links
    % urlcolor=blue,% color of external links
    hidelinks,% hide links color and box
}
%---------------------------------------------------------------------------%
%->> END
%---------------------------------------------------------------------------%
%---------------------------------------------------------------------------%
%- 辅助命令，后文中的所有\include均可在此单独列出，用逗号隔开，
%- 以此只编译必要的章节，加快编译速度，待全文完成后可注释本命令，即可编译全文。
%- 也可注释部分内容，正文中所有的内容均可注释后避免其参与编译，包含maketitle等命令，
%- 执行此命令或注释后可能导致章节序号发生错误，无需担心，全文编译后即可恢复
% \includeonly{Tex/Abstract,Tex/Appendix}
% %---------------------------------------------------------------------------%
\begin{document}
%-
\maketitle% 制作封面
%-
%- 声明包含两种形式，
%- 如果参数为空则可自动生成默认声明页，也可设置参数导入签字后的扫描版PDF文件
% \makedeclaration{declaration}% 制作声明，参数为扫描版文件名，默认在Image下
\makedeclaration{}% 制作声明，自动生成
%-
\frontmatter% 初始化摘要页环境，不建议注释
%-
%---------------------------------------------------------------------------%
%->> Abstract
%---------------------------------------------------------------------------%
%-
%-> 中文摘要
%-
\begin{abstract}
    本文是深圳大学学位论文模板 szuthesis 的使用说明文档。
    主要内容为介绍 \LaTeX{} 文档类 szuthesis 的用法，以及如何使用 \LaTeX{} 快速高效地撰写学位论文。

    \keywords{深圳大学，学位论文，\LaTeX{}模板}% 中文关键词
\end{abstract}
%-
%-> 英文摘要
%-
\begin{ABSTRACT}
    This paper is a help documentation for the \LaTeX{} class szuthesis,
    which is a thesis template for the Shenzhen University.
    The main content is about how to use the szuthesis, as well as how to write thesis
    efficiently by using \LaTeX{}.

    \KEYWORDS{Shenzhen University, Thesis, \LaTeX{} Template}% 英文关键词
\end{ABSTRACT}
%---------------------------------------------------------------------------%
%-
\tableofcontents% 目录
%-
\mainmatter% 初始化正文环境，不建议注释
%-
\chapter{引言}\label{chap:intro}

\section{研究背景}\label{sec:background}

深圳大学学位论文模板 szuthesis 基于中国科学院大学学位论文模板 ucasthesis\footnote{\url{https://github.com/mohuangrui/ucasthesis}} 二次开发。
相对于 ucasthesis ，本项目进行了许多改动，关键的变化如下：

\begin{enumerate}[wide=\parindent]
    \item 修改相关配置，使其符合深圳大学学位论文印刷格式样式；
    \item 封装了更多的命令，极大化简了 Thesis.tex 中的内容；
    \item 化简了许多操作，例如化简了字体相关补丁，遵循ctex字体关于粗体和斜体的相关设定，化简了参考文献格式等；
    \item 参考了上海交通大学论文模板 SJTUThesis\footnote{\url{https://github.com/sjtug/SJTUThesis}} 的设计思路。
\end{enumerate}

考虑到许多同学可能缺乏 \LaTeX{} 使用经验，szuthesis 不仅延续了 ucasthesis 的轻易使用的特点，同时封装了更多功能，极大简化了相关逻辑。
在本文中，对用 \LaTeX{} 撰写论文的一些主要难题，如制图、制表、文献索引等，进行了详细说明，并提供了相应的代码样本，大大降低了初学者的学习成本。
所以，如果你是初学者，请不要放弃，因为同样为初学者的我，十分明白让 \LaTeX{} 简单易用的重要性，而这正是 szuthesis 所追求和体现的。

当前 szuthesis 模板完美遵循《学术学位硕士学位论文印刷格式样式》与《专业硕士学位论文印刷格式样式》中规定的学位论文撰写要求和封面设定。
兼顾操作系统 Windows、Linux、MacOS；目前仅支持 Xe\LaTeX{} 引擎；文献编译引擎biber (biblatex)。
支持中文书签、中文渲染、拷贝 PDF 中的文本到其他文本编辑器等特性。

szuthesis 的目标在于简化学位论文的撰写，充分利用 \LaTeX{} 的优点，例如，使用 git 进行版本控制，专注写作且无需关注格式，更便利的公式书写环境，支持注释等。


\section{系统要求}\label{sec:system}

szuthesis 宏包可以在目前主流的 \href{https://en.wikibooks.org/wiki/LaTeX/Introduction}{\LaTeX{}} 发行版中使用，
如 \TeX{}Live 和 MiK\TeX{}。因 C\TeX{} 套装已停止维护，\textbf{请勿使用}。
请勿混淆 C\TeX{} 套装\footnote{\url{http://www.ctex.org/CTeX}}与 C\TeX{} 宏集\footnote{\url{https://ctan.org/pkg/ctex?lang=en}}。
C\TeX{} 套装基于 Windows 下的 MiKTeX 开发，在其基础上增加了对中文的完整支持，已于 2012 年起停止维护。
而 C\TeX{} 宏集是通用 \LaTeX{} 排版框架，为中文 \LaTeX{} 文档提供了汉字支持，主要包括宏包 ctex 以及中文文档类 ctexart、 ctexbook 等。

推荐的 \LaTeX{} 发行版如下：

\begin{center}
    %\footnotesize% fontsize
    %\setlength{\tabcolsep}{4pt}% column separation
    %\renewcommand{\arraystretch}{1.5}% row space 
    \begin{tabular}{lc}
        \toprule
        操作系统         & \LaTeX{}发行版                                                                                        \\
        \midrule
        Linux 或 Windows & \href{https://www.tug.org/texlive/}{\TeX{}Live Full} 或 \href{https://miktex.org/download}{MiK\TeX{}} \\
        MacOS            & \href{https://www.tug.org/mactex/}{Mac\TeX{} Full} 或 \href{https://miktex.org/download}{MiK\TeX{}}   \\
        \bottomrule
    \end{tabular}
\end{center}

请从各软件官网下载安装程序，勿使用不明程序源。若系统原带有旧版的 \LaTeX{} 发行版并想安装新版，请\textbf{先完全卸载旧版再安装新版}。
推荐安装2019年后的版本。可能因为网络问题导致安装速度较慢，推荐在安装时无需选择额外宏包，
安装完成后添加清华源\footnote{\url{https://mirrors.tuna.tsinghua.edu.cn/help/CTAN/}}，再继续安装所需宏包。

如选择部分安装，请安装后检测以下宏包知否安装，若未安装导致的BUG不易排查：

\begin{center}
    \small% fontsize
    \renewcommand{\arraystretch}{0.8}% row space 
    \begin{tabular}{ll}
        \toprule
        宏包                 & 功能                    \\
        \midrule
        xits                 & 开源Times New Roman字体 \\
        biber                & biblatex引擎            \\
        biblatex-gb7714-2015 & biblatex格式            \\
        latexmk              & 自动编译latex文档       \\
        \bottomrule
    \end{tabular}
\end{center}

安装 \LaTeX{} 发行版后，即可使用任意编辑器开始书写。
在这里推荐使用 VS Code\footnote{\url{https://code.visualstudio.com/}} 作为编辑器。一方面其可单纯的作为编辑器使用，
同时又可以搭配插件进行扩展。可以搭配 Git 进行版本控制，又可以安装 LaTeX Workshop 插件直接进行编译。
LaTeX Workshop 插件提供了诸如 Linting，Formatting，Intellisense，PDF 文件预览，公式预览，
全文大纲等诸多功能\footnote{\url{https://github.com/James-Yu/LaTeX-Workshop\#features-taster}}。
使用快捷键可以极大提高编写效率，例如使用 \lstinline!Ctrl+Alt+j! 可以快速从 tex 文本跳转到 PDF 中对应的位置，
而在PDF预览中使用 \lstinline!Ctrl+鼠标左键! 就可以快速定位对应的 tex 文本。


\section{编译}

\begin{enumerate}[wide=\parindent]
    \item 安装软件：根据所用操作系统和章节~\ref{sec:system} 中的信息安装 \LaTeX{} 编译环境。

    \item 获取模板：下载 szuthesis 项目。szuthesis 不仅提供了相应的模板，同时也提供了编译样例，
          下载时推荐下载整个 szuthesis 文件夹，而不是单独的 cls 文档类。

    \item 编译模板：参考项目主页编译部分。
\end{enumerate}

编译完成后即可获得这份说明文档。而这也完成了学习使用 szuthesis 撰写论文的一半进程。
什么？这就学成一半了，这么简单？？？是的，就这么简单！


\section{文档目录简介}

\subsection{Thesis.tex}

Thesis.tex 为主文档，包含了论文全篇的主要架构。其中，document 中的所有内容均可注释后避免其参与编译，包含 maketitle 等命令，
注释后可能导致章节序号发生错误，无需担心，全文编译后即可恢复。注释后可加快编译速度，例如参考文献页无须随文档实时编译，
只需要全文完成后编译参考文献页即可，这也是使用 \LaTeX{} 编写文档的优点之一。

\subsection{Temp文件夹}

编译后，生成的临时文件皆存于Temp文件夹内，包括编译得到的 PDF 文档，其存在是为了保持工作空间的整洁，因为好的心情是很重要的。

\subsection{szuthesis.cls}

\verb!texmf\tex\latex\szuthesis\szuthesis.cls! 目录下的 szuthesis.cls 为文档类，定义了论文的核心格式，
包括论文排版、引用格式、页眉页脚、字体字号等。其中，根据《印刷格式样式》规定，参考文献后的字号均与正文字号不同。

\subsection{config.tex}

szuthesis.cls 需要传入一些参数用来生成封面信息，config.tex 可用来传入这些参数。
后边则定义了一些可选的宏包，这些宏包并不完全属于《印刷格式样式》规定的排版，可以自由选择是否启用。
例如数学公式的字体、代码片段、超链接等，均在 config.tex 进行了定义，这些可以根据需要对它进行调整。

\subsection{Tex文件夹}

文件夹内为论文的所有正文内容，这也是使用 szuthesis 撰写学位论文时，主要关注和修改的一个位置。
\textbf{注：所有文件都必须采用 UTF-8 编码，否则编译后将出现乱码文本}，详细分类介绍如下：

\begin{itemize}
    \item Abstract.tex：摘要信息。
    \item ChapterX.tex：论文的各个章节，可根据需要添加和撰写。\textbf{添加新章节时，注意编码格式，可拷贝一个已有的章文件再重命名，以继承文档的 UTF-8 编码}。
    \item Appendix.tex：附录，注意附录字号与正文不同，仅用于添加补充信息，如有整段文本建议放置于正文中。
    \item Acknowledgements.tex：致谢。
    \item Publications.tex: 研究成果。
\end{itemize}

\subsection{Image文件夹}

用于放置论文中所需要的图形类文件，支持格式有：jpg, png, pdf 等，需要更多支持格式可在 config.tex 中配置。
不建议为各章节图片建子目录，即使图片众多，若命名规则合理，图片查询亦是十分方便。

\subsection{Biblio文件夹}

ref.bib 为参考文献信息，可在 config.tex 中配置。

\subsection{.vscode文件夹}

这一文件夹用于保存 VS Code 的配置文件，其中 settings.json 保存了部分 latexmk 所需的配置项。



\section{帮助与问题反馈}\label{sec:help}

对于 \LaTeX{} 相关问题，推荐使用 texdoc 命令查阅相关文档。
例如安装 lshort-chinese 宏包后，可使用 \textbf{texdoc lshort-chinese} 命令打开一份教程，包含了 \LaTeX{} 入门相关的知识。
使用 \textbf{texdoc ctex} 则可打开 ctex 宏包的文档，包含中文排版相关的内容，例如第5节中则详细介绍了中文字体字号。
大多数宏包都提供了非常详尽的文档，都可以使用 texdoc 查阅。

欢迎各位同学提出宝贵意见，一起不断改进模板。

\chapter{\LaTeX{} 入门}

\section{二级标题}

\subsection{三级标题}

\subsubsection{四级标题}

深圳大学1983年经国家教育部批准设立。中央、教育部和地方高度重视特区大学建设。北大援建中
文、外语类学科，清华援建电子、建筑类学科，人大援建经济、法律类学科，一大批知名学者云集
深圳大学。建校伊始，学校在高校管理体制上锐意改革，在奖学金、学分制、勤工俭学等方面进行
了积极探索，率先在国内实行毕业生不包分配和双向选择制度，推行教职员工全员聘任制度和后勤
部门社会化管理改革，在全国引起强烈反响\footnote{这是脚注}。

\section{字体}

字体部分主要分为正文字体和数学公式字体，这里主要展示正文字体，数学公式字体详见下节。

正文字体延续 ctex 关于字体的相关设置，需要注意的是，中文字体通常无粗体与斜体，例如
Windows 平台下，正文默认宋体，粗体则使用黑体代替，斜体则使用了楷体。MacOS 对多数
字体提供了完整的粗体表示，但是默认的斜体依然使用了楷体。通常论文的正文应简洁明了，
尽量避免使用多种字体。如需使用粗体或斜体，需重新定义默认字体族或者定义新的字体族。

由上段可知，字体的选择依赖于字库，而字库又依赖于系统设置，模板支持windows、adobe、mac、ubuntu
四种中文字库，通常情况下可根据系统类型自动选择。但是 Ubuntu 系统因缺少仿宋字体，而部分章节的标题
会使用到仿宋，模板使用宋体进行了代替，依然推荐 Ubuntu 用户安装完整 adobe 字体并且指定字库。

关于字库的显示效果可参考表~\ref{tab:fontset}，该表展示了 Windows 系统下主要的字体以及对应的粗体与楷体的
兼容情况。可以看到默认的粗体使用了黑体代替，而斜体则使用了楷书。中文字体的设置均不影响英文字体的设置，英文
字体除摘要页标题使用了 Lucida Console 字体外，全局使用开源 Times New Roman 字体 XITS。可以看到 \cs{textit}
和 \cs{kaishu} 都可以生成中文的楷书，但是对于英文字体却产生了不同的影响，书写时需注意。

\def\SZUFONTEG{深圳大学fgFG456}
\begin{table}
  \centering
  \caption{字体选择}\label{tab:fontset}
  \begin{tabular}{lcccc}
    \toprule
    字体 & 命令          & 常规                  & 粗体                         & 斜体                         \\
    \midrule
    宋体 & -             & {\SZUFONTEG}          & \textbf{\SZUFONTEG}          & \textit{\SZUFONTEG}          \\
         & \cs{textrm}   & \textrm{\SZUFONTEG}   & \textbf{\textrm{\SZUFONTEG}} & \textit{\textrm{\SZUFONTEG}} \\
         & \cs{songti}   & {\songti\SZUFONTEG}   & -                            & -                            \\
    黑体 & \cs{textsf}   & \textsf{\SZUFONTEG}   & \textbf{\textsf{\SZUFONTEG}} & -                            \\
         & \cs{heiti}    & {\heiti\SZUFONTEG}    & -                            & -                            \\
    仿宋 & \cs{texttt}   & \texttt{\SZUFONTEG}   & -                            & -                            \\
         & \cs{fangsong} & {\fangsong\SZUFONTEG} & -                            & -                            \\
    楷书 & \cs{textit}   & \textit{\SZUFONTEG}   & -                            & -                            \\
         & \cs{kaishu}   & {\kaishu\SZUFONTEG}   & -                            & -                            \\
    \bottomrule
  \end{tabular}
\end{table}

此外，windows 和 mac 对应的中文字库还提供了微软雅黑 \cs{yahei}、隶书 \cs{lishu}、幼圆 \cs{youyuan} 三种字体，
\cs{yahei} 默认提供了粗体。更多关于默认字体与字体库的选择相关问题可查阅 ctex 文档 4.3 节中文字库。

除了中文字库外，由 XITS 字库可支持俄语 Привет ，默认的中文字体宋体通常支持日语 こんにちは 。
如需要韩语，则需要导入新的字库，如 Windows 通常带有 ARIALUNI.TTF 韩语字库，可以使用
\lstinline[language=Tex]!\newCJKfontfamily{\hanyu}{ARIALUNI.TTF}! 导入新的字库，就可以使用 \cs{hanyu} 来
书写韩语了。其他语言同理，可以使用这个命令导入对应的字库。

\section{数学}

\subsection{数学符号与公式}

数学公式均统一使用了开源 Times New Roman 字体 XITS MATH，这套字体除了基本的数字与字母外，还额外设计了许多
数学符号，包括全局粗体。使用这套字体需导入 \pkg{unicode-math}，所以在启用数学相关字体时需要注意，使用 \cs{sym..} 系列
命令代替 \cs{math..} 系列命令，以获得统一的字体风格，\cs{math..} 系列命令将会使用普通字母的字体 XITS，而不是
对应的 XITS MATH。例如加粗 \cs{mathbf} 将使用 \cs{symbf} 代替。
\[
  \begin{split}
    \symcal{L} \{f\}(s) &= \int _{0^{-}}^{\infty} f(t) e^{-st} \, \symrm{d}t, \
    \symscr{L} \{f\}(s) = \int _{0^{-}}^{\infty} f(t) e^{-st} \, \symrm{d}t\\
    \symcal{F} \bigl( f(x+x_{0}) \bigr) &= \symcal{F} \bigl( f(x) \bigr) e^{2\pi i\xi x_{0}}, \
    \symscr{F} \bigl( f(x+x_{0}) \bigr) = \symscr{F} \bigl( f(x) \bigr) e^{2\pi i\xi x_{0}}
  \end{split}
\]
\begin{equation}\label{equ:1}
  \begin{cases}
    \frac{\partial \rho}{\partial t} + \nabla\cdot(\rho\Vector{V}) = 0                                                \\
    \frac{\partial (\rho\Vector{V})}{\partial t} + \nabla\cdot(\rho\Vector{V}\Vector{V}) = \nabla\cdot\Tensor{\sigma} \\
    \frac{\partial (\rho E)}{\partial t} + \nabla\cdot(\rho E\Vector{V}) = \nabla\cdot(k\nabla T) + \nabla\cdot(\Tensor{\sigma}\cdot\Vector{V})
  \end{cases}
\end{equation}
\begin{equation}\label{equ:2}
  \frac{\partial }{\partial t}\int\limits_{\Omega} u \, \symrm{d}\Omega + \int\limits_{S} \Vector{n}\cdot(u\Vector{V}) \, \symrm{d}S = \dot{\phi}
\end{equation}

表~\ref{tab:mathfontset} 中简要示范了几种常用的数学字体，可以看到指令并不是对所有的数字和希腊字母都有效。更多关于
数学字体相关的知识推荐查阅《在 \LaTeX{} 中使用 OpenType 字体》\footnote{\url{https://stone-zeng.github.io/2020-05-02-use-opentype-fonts-iii/}\label{stonezeng}}
一文，该文中``字符支持''一节详细介绍了各个命令的使用场景。

\def\SZUMATHFONTEG{fgFG456\delta\Delta}
\begin{table}[!htbp]
  \centering
  \caption{数学公式字体选择}\label{tab:mathfontset}
  \begin{tabular}{cccccc}
    \toprule
    命令        & 样例                        & 命令         & 样例                         & 命令         & 样例                         \\
    \midrule
    -           & \({\SZUMATHFONTEG}\)        & \cs{symbf}   & \(\symbf{\SZUMATHFONTEG}\)   & \cs{symrm}   & \(\symrm{\SZUMATHFONTEG}\)   \\
    \cs{symup}  & \(\symup{\SZUMATHFONTEG}\)  & \cs{symbfup} & \(\symbfup{\SZUMATHFONTEG}\) & \cs{symit}   & \(\symit{\SZUMATHFONTEG}\)   \\
    \cs{symbb}  & \(\symbb{\SZUMATHFONTEG}\)  & \cs{symsfup} & \(\symsfup{\SZUMATHFONTEG}\) & \cs{symtt}   & \(\symtt{\SZUMATHFONTEG}\)   \\
    \cs{symcal} & \(\symcal{\SZUMATHFONTEG}\) & \cs{symscr}  & \(\symscr{\SZUMATHFONTEG}\)  & \cs{symfrak} & \(\symfrak{\SZUMATHFONTEG}\) \\
    \bottomrule
  \end{tabular}
\end{table}

在使用数学字体过程中，应遵循相应的标准，根据 GB 3102.11《物理科学和技术中使用的数学
符号》\footnote{\url{https://zh.wikisource.org/wiki/GB_3102.11}}中规定：

\begin{itemize}
  \small
  \item 变量 (例如 \(x\), \(y\) 等)、变动附标 (例如 \(\sum_i x_i\) 中的 \(i\)) 及函数 (例如 \(f\), \(g\) 等)
        用斜体字母表示。点 \(A\)、线段 \(AB\) 及弧 \(CD\) 用斜体字母表示。
        在特定场合中视为常数的参数 (例如 \(a\), \(b\) 等) 也用斜体字母表示。
  \item 有定义的已知函数 (例如 \(\sin \)，\(\exp \)，\(\ln \)，\(\Gamma \) 等) 用正体字母表示。
  其值不变的数学常数 (例如 \(\exp = 2.718 281 8\ldots \)，\(\symrm{\pi} = 3.141 592 6\ldots \)，\(\symrm{i}^2 = −1\) 等) 用正体
  字母表示。已定义的算子 (例如 \(\Div \)，\(\deltarm x\) 中的 \(\deltarm \) 及 \(\Dif 𝑓/\Dif 𝑥\) 中的 \(\Dif \)) 也用正体字母表示。
  \item 数字表中数 (例如 351 204，1.32，7/8) 的表示用正体。
\end{itemize}

根据这一规定，在 config.tex 中额外定义了一些通用命令，如表~\ref{tab:mathcmd} 所示。使用这些命令，
我们可以方便的定义公式，如公式\ref{equ:1}、\ref{equ:2}和公式~\ref{equ:3}\textsuperscript{\ref{stonezeng}}。
\begin{equation}\label{equ:3}
  \begin{aligned}
    \oiint_{\partial\Omega} \Vector{}{E} \cdot \Dif \Vector{S} & = \frac{1}{\epsilon_0} \iiint_\Omega \rho \, \Dif V & \quad &
    \oint_{\partial\Sigma} \Vector{E} \cdot \Dif \Vector{l} = -\frac{\Dif}{\Dif t} \iint_\Sigma \Vector{B} \cdot \Dif \Vector{S}         \\
    \oiint_{\partial\Omega} \Vector{B} \cdot \Dif \Vector{S} & = 0                                                & \quad &
    \oint_{\partial\Sigma} \Vector{B} \cdot \Dif \Vector{l} = \mu_0 \iint_\Sigma \Vector{J} \cdot \Vector{S}
    + \mu_0 \epsilon_0 \frac{\Dif}{\Dif t} \iint_\Sigma \Vector{E} \cdot \Dif \Vector{S} \\
  \end{aligned}
\end{equation}

\begin{table}[!htbp]
  \centering
  \caption{自定义指令}\label{tab:mathcmd}
  \begin{tabular}{cccccccc}
    \toprule
    用法                   & 样例            & 用法                   & 样例           & 用法                   & 样例           & 用法       & 样例        \\
    \midrule
    \verb!\Vector{A}! & \(\Vector{A}\)  & \verb!\Matrix{A}! & \(\Matrix{A}\) & \verb!\Tensor{A}! & \(\Tensor{A}\) & \cs{Dif x} & \(\Dif x\)  \\
    \verb!\Const{\pi}! & \(\Const{\pi}\) & \cs{deltarm}           & \(\deltarm \)  & \cs{Div}               & \(\Div \)      & \cs{Trace} & \(\Trace \) \\
    \bottomrule
  \end{tabular}
\end{table}

\subsection{标题中带有公式 \texorpdfstring{\(\Lambda,\lambda,\theta,\bar{\Lambda},\sqrt{S_{NN}}\)}{}}

\LaTeX{} 支持在标题中带有公式，以及生成的目录中带有公式。如本节标题所示，在目录中均可显示完整的公式。
由于 \pkg{hyperref} 在生成书签时不支持公式，所以当使用该包时，可以使用 \cs{texorpdfstring} 命令。
这个命令的第一个参数可以书写公式，第二个参数则用于生成书签，可以使用 \cs{textLambda} 等书签兼容的公式，
不过最简单的方法就是空出第二个参数，因为书签在打印时不显示，所以无需额外关注。

\subsection{数学环境}
以下环境共享一个计数器，如果取消这一设定，在 config.tex 中删除 \cs{newtheorem} 的第二个参数[theorem]即可。

\begin{theorem}      这是一个定理。 \end{theorem}
\begin{axiom}        这是一个公理。 \end{axiom}
\begin{lemma}        这是一个引理。 \end{lemma}
\begin{corollary}    这是一个推论。 \end{corollary}
\begin{assertion}    这是一个断言。 \end{assertion}
\begin{proposition}  这是一个命题。 \end{proposition}
\begin{conjecture}   这是一个猜想。 \end{conjecture}
\begin{assumption}   这是一个假设。 \end{assumption}

以下环境均为独立的计数器。

\begin{definition}   这是一个定义。 \end{definition}
\begin{example}      这是一个例子。 \end{example}
\begin{problem}      这是一个问题。 \end{problem}
\begin{exercise}     这是一个练习。 \end{exercise}
\begin{remark}       这是一个注。   \end{remark}
\begin{solution}     这是一个解。   \end{solution}
\begin{proof}        这是一个证明。 \end{proof}

\section{浮动体}

\subsection{图形}

插入单个图形时效果如图~\ref{fig:single}。同样的方法可以插入多个图形，通过可选参数可以指定图形的高度
或者宽度，效果如图~\ref{fig:multi}。如果各个子图相互独立，可以使用 \pkg{minipage}，如图~\ref{fig:parallel1}
和图~\ref{fig:parallel2}。

\begin{figure}[!htbp]
  \centering
  \includegraphics{logo} \\
  此处可撰写对图的一些说明 \\
  支持多行
  \caption{此处是图的标题，也可以在此处撰写说明}\label{fig:single}
\end{figure}

\begin{figure}[!htbp]
  \centering
  \includegraphics[height=2cm]{logo}
  \hspace{1cm}
  \includegraphics[height=2cm]{logo}
  \caption{插入多个图形，共享同一个计数器}\label{fig:multi}
\end{figure}

\begin{figure}[!htbp]
  \begin{minipage}{0.45\textwidth}
    \centering
    \includegraphics[height=2cm]{logo}
    \caption{并排图一}\label{fig:parallel1}
  \end{minipage}
  \hfill
  \begin{minipage}{0.45\textwidth}
    \centering
    \includegraphics[height=2cm]{logo}
    \caption{并排图二}\label{fig:parallel2}
  \end{minipage}
\end{figure}

如果要为共用一个计数器的多个子图添加子标题，建议使用较新的 \pkg{subcaption} 宏包，
不建议使用 \pkg{subfigure} 或 \pkg{subfig} 等宏包。可以直接使用 \cs{subcaptionbox}
并排子图，或者使用 \cs{subcaption} 直接放在 minipage 中，用法同 \cs{caption}。
如图~\ref{fig:subcapbox}，通过对 \cs{subcaptionbox} 和 \cs{includegraphics} 设置不同
的尺寸，即可实现两幅子图的任意的比例和尺寸。

\begin{figure}[!htbp]
  \centering
  \subcaptionbox{图一标题}[0.6\textwidth]{\includegraphics[height=2cm]{logo}}
  \subcaptionbox{图二标题}[0.3\textwidth]{\includegraphics[height=1cm]{logo}}
  \caption{同一个计数器包含多个子图}\label{fig:subcapbox}
\end{figure}

\subsection{表格}

根据印刷格式样式，表格中应使用小五号宋体，最好能够同时调整行间距和表格间距以获得较好的效果。
学术论文通常使用 \cs{toprule}、\cs{midrule}、\cs{bottomrule} 绘制三线表，效果如表~\ref{tab:sample} 所示。
此外，\pkg{threeparttable} 宏包提供了子环境可以支持在表格中添加脚注，\pkg{longtable} 宏包实现了可以书写跨页
的长表格。

\begin{table}[!htbp]
  \caption{这是一个标准的三线表}\label{tab:sample}
  \centering
  \footnotesize% fontsize，小五号字
  \setlength{\tabcolsep}{4pt}% column separation，表格间距
  \renewcommand{\arraystretch}{1.2}% row space ，行间距
  \begin{tabular}{lcccccccc}
      \toprule
      行号 & \multicolumn{8}{c}{跨多列的标题}\\
      \midrule
      Row 1 & $1$ & $2$ & $3$ & $4$ & $5$ & $6$ & $7$ & $8$\\
      Row 2 & $1$ & $2$ & $3$ & $4$ & $5$ & $6$ & $7$ & $8$\\
      Row 3 & $1$ & $2$ & $3$ & $4$ & $5$ & $6$ & $7$ & $8$\\
      \bottomrule
  \end{tabular}
\end{table}

\section{代码}

在论文中插入代码需要 \pkg{listings} 宏包，它支持代码的语法高亮等特性。config.tex 中配置了最简单的代码风格，
没有颜色，可以用于打印。

lstlisting 环境提供了可选参数，例如 captain 或者 title，前者可显示序号，而后者只显示标题而没有序号。
可选参数 label 提供了交叉引用的功能，例如代码~\ref{code}。

\begin{lstlisting}[language=C,caption=代码样例,label=code]
#include <stdio.h>
#include <unistd.h>
#include <sys/types.h>
#include <sys/wait.h>

int main() {
  pid_t pid;

  switch ((pid = fork())) {
  case -1:
    printf("fork failed\n");
    break;
  case 0:
    /* child calls exec */
    execl("/bin/ls", "ls", "-l", (char*)0);
    printf("execl failed\n");
    break;
  default:
    /* parent uses wait to suspend execution until child finishes */
    wait((int*)0);
    printf("is completed\n");
    break;
  }

  return 0;
}\end{lstlisting}

在 tex 文件中书写代码需要格外注意缩进，更加方便地是直接从代码源文件中导入代码片段，
即使用 \cs{lstinputlisting} 命令，它提供了参数 linerange 可以指定导入源文件的哪几行，
通过 \cs{lstname} 还可以直接获取文件名。

\lstinputlisting[linerange={10-17},language=C,caption=\lstname]{Source/code.cpp}

书写算法可以使用 \pkg{algorithm2e}，如算法~\ref{algo}，更多相关内容可查询相关文档。

\begin{algorithm}[!htbp]
    \SetAlgoLined\normalem\linespread{1.2}\small
    \caption{How to write algorithms}\label{algo}
    \KwData{this text}
    \KwResult{how to write algorithm with \LaTeX2e }
    initialization\;
    \While{not at end of this document}{
        read current\;
        \eIf{understand}{
            go to next section\;
            current section becomes this one\;
        }{
            go back to the beginning of current section\;
        }
    }
\end{algorithm}

\section{参考文献}

参考文献应符合国标 GB/T 7714-2015，该标准为替换印刷格式样式中规定的国标 GB/T 7714-2005 的最新版本，
使用 Bib\LaTeX{} 以及样式包 \pkg{biblatex-7714-2015} 可以轻松实现，后端使用 biber 引擎。

通常可以使用 \cs{cite} 直接引用，例如\cite{lamport1986document}，参数可以添加多个，用逗号隔开，可以实现
同时引用多个参考文献\cite{chen2005zhulu,chu2004tushu,stamerjohanns2009mathml}。
使用 \cs{parencite} 可以获得非上标的以用\parencite{betts2005aging}，同样可以同时引用
多个\parencite{bravo1990comparative,hls2012jinji,niu2013zonghe}。此外，还支持在脚注的引用，
这种方式的引用将与脚注共享编号，并且同时也会打印在最后的参考引用页中\footfullcite{chen1980zhongguo}，打印在
参考引用页中的序号与脚注编号无关。

参考文献页\cite{yuan2012lana,yuan2012lanb,yuan2012lanc}将按照引用出现的顺序\footfullcite{walls2013drought}
生成序号\parencite{Bohan1928,Dubrovin1906}。

% \include{Tex/Chapter3}
%- 
\backmatter% 初始化其他部分环境，不建议注释
%-
\szubibliography% 导入参考文献
%-
\chapter[附录]{附\quad{}录}

\section*{论文无附录者无需附录部分}

附录中的标题尽量使用 \cs{section*} 的形式。

\section*{测试生僻字}

霜蟾盥薇曜灵霜颸妙鬘虚霩淩澌菀枯菡萏泬寥窅冥毰毸濩落霅霅便嬛岧峣瀺灂姽婳愔嫕飒纚棽俪緸冤
莩甲摛藻卮言倥侗椒觞期颐夜阑彬蔚倥偬澄廓簪缨陟遐迤逦缥缃鹣鲽憯懔闺闼璀错媕婀噌吰澒洞阛闠
覼缕玓瓑逡巡諓諓琭琭瀌瀌踽踽叆叇氤氲瓠犀流眄蹀躞赟嬛茕頔璎珞螓首蘅皋惏悷缱绻昶皴皱颟顸愀
然菡萏卑陬纯懿犇麤掱暒墌墍墎墏墐墒墒墓墔墕墖墘墖墚墛坠墝增墠墡墢墣墤墥墦墧墨墩墪樽墬墭堕
墯墰墱墲坟墴墵垯墷墸墹墺墙墼墽垦墿壀壁壂壃壄壅壆坛壈壉壊垱壌壍埙壏壐壑壒压壔壕壖壗垒圹垆
壛壜壝垄壠壡坜壣壤壥壦壧壨坝塆圭嫶嫷嫸嫹嫺娴嫼嫽嫾婳妫嬁嬂嬃嬄嬅嬆嬇娆嬉嬊娇嬍嬎嬏嬐嬑嬒
嬓嬔嬕嬖嬗嬘嫱嬚嬛嬜嬞嬟嬠嫒嬢嬣嬥嬦嬧嬨嬩嫔嬫嬬奶嬬嬮嬯婴嬱嬲嬳嬴嬵嬶嬷婶嬹嬺嬻嬼嬽嬾嬿
孀孁孂娘孄孅孆孇孆孈孉孊娈孋孊孍孎孏嫫婿媚嵭嵮嵯嵰嵱嵲嵳嵴嵵嵶嵷嵸嵹嵺嵻嵼嵽嵾嵿嶀嵝嶂嶃
崭嶅嶆岖嶈嶉嶊嶋嶌嶍嶎嶏嶐嶑嶒嶓嵚嶕嶖嶘嶙嶚嶛嶜嶝嶞嶟峤嶡峣嶣嶤嶥嶦峄峃嶩嶪嶫嶬嶭崄嶯嶰
嶱嶲嶳岙嶵嶶嶷嵘嶹岭嶻屿岳帋巀巁巂巃巄巅巆巇巈巉巊岿巌巍巎巏巐巑峦巓巅巕岩巗巘巙巚帠帡帢
帣帤帨帩帪帬帯帰帱帲帴帵帷帹帺帻帼帽帾帿幁幂帏幄幅幆幇幈幉幊幋幌幍幎幏幐幑幒幓幖幙幚幛幜
幝幞帜幠幡幢幤幥幦幧幨幩幪幭幮幯幰幱庍庎庑庖庘庛庝庠庡庢庣庤庥庨庩庪庬庮庯庰庱庲庳庴庵庹
庺庻庼庽庿廀厕廃厩廅廆廇廋廌廍庼廏廐廑廒廔廕廖廗廘廙廛廜廞庑廤廥廦廧廨廭廮廯廰痈廲廵廸廹
廻廼廽廿弁弅弆弇弉弖弙弚弜弝弞弡弢弣弤弨弩弪弫弬弭弮弰弲弪弴弶弸弻弼弽弿彖彗彘彚彛彜彝彞
彟彴彵彶彷彸役彺彻彽彾佛徂徃徆徇徉后徍徎徏径徒従徔徕徖徙徚徛徜徝从徟徕御徢徣徤徥徦徧徨复
循徫旁徭微徯徰徱徲徳徴徵徶德徸彻徺忁忂惔愔忇忈忉忔忕忖忚忛応忝忞忟忪挣挦挧挨挩挪挫挬挭挮
挰掇授掉掊掋掍掎掐掑排掓掔掕挜掚挂掜掝掞掟掠采探掣掤掦措掫掬掭掮掯掰掱掲掳掴掵掶掸掹掺掻
掼掽掾掿拣揁揂揃揅揄揆揇揈揉揊揋揌揍揎揑揓揔揕揖揗揘揙揤揥揦揧揨揫捂揰揱揲揳援揵揶揷揸揻
揼揾揿搀搁搂搃搄搅搇搈搉搊搋搌搎搏搐搑搒摓摔摕摖摗摙摚摛掼摝摞摠摡斫斩斮斱斲斳斴斵斶斸旪
旫旮旯晒晓晔晕晖晗晘晙晛晜晞晟晠晡晰晣晤晥晦晧晪晫晬晭晰晱晲晳晴晵晷晸晹晻晼晽晾晿暀暁暂
暃暄暅暆暇晕晖暊暋暌暍暎暏暐暑暒暓暔暕暖暗旸暙暚暛暜暝暞暟暠暡暣暤暥暦暧暨暩暪暬暭暮暯暰
昵暲暳暴暵
% 导入附录
%-
%- 2020年新增，添加答辩记录，建议分成三个独立的PDF文件，
%- \szuaddpdf命令包含两个参数，[]中为可选参数，用于生成目录，{}中为PDF文件名，默认在Image下
% \szuaddpdf[指导教师对研究生学位论文的学术评语]{pingyu}
% \szuaddpdf[学位论文答辩委员会决议书]{dabian1}
%-\szuaddpdf{dabian2}% 前一页生成目录即可
%-
\chapter[致谢]{致\quad{}谢}

感谢国科大学位论文ucasthesis项目，以及上海交大学位论文SJTUThesis项目对开源的无私贡献。

感谢CTeX-kit项目对\LaTeX{}提供了中文支持。

本项目离不开\LaTeX{}开源社区的诸多宝贵资料，对\LaTeX{}开源社区各位贡献者表示感谢。

感谢参与测试和使用的各位同学。% 导入致谢
%-
\chapter{攻读硕士学位期间的研究成果}

%- 可以直接使用引用的格式

\begin{enumerate}[label = {[\arabic*]}]
    \item TANG T T. xxx xxxx xx xxx[M]. IEEE Transactions on xxxx. 2020-01-01.
    \item 唐同学. 一种基于某某的方法:中国,01100000.5[P]. 2020-01-01.
\end{enumerate}


%- 或者区分开

% \section*{论文}

% \begin{enumerate}[label = {[\arabic*]}]
%     \item TANG T T. xxx xxxx xx xxx. IEEE Transactions on xxxx. 2020.
% \end{enumerate}

% \section*{专利}

% \begin{enumerate}[label = {[\arabic*]}]
%     \item 唐同学. 一种基于某某的方法:中国,01100000.5. 2020-01-01.
% \end{enumerate}% 导入研究成果
%-
\end{document}
% %---------------------------------------------------------------------------%
